\subsection*{幂运算先需要一个定义}
实数正底数的幂可以写成 $a^b=\exp(b\ln a)$。把底数扩展到复数，首先要解决复对数的问题。复指数绕整圈回到同一点，因此 $\mathrm i$ 对应角度 $\pi/2+2\pi k$，而非只有一个角度。

选择其中一个对数值 $\log\mathrm i=\mathrm i(\pi/2+2\pi k)$，再乘以指数 $\mathrm i$，得到实数 $-(\pi/2+2\pi k)$。对它取指数，结果便是正实数。所谓“虚数的虚数次方是实数”，关键并非神秘抵消，而是两次乘 $\mathrm i$ 带来负号。

\subsection*{主值是约定，多值是结构}
选择主对数对应 $k=0$，得到 $\mathrm e^{-\pi/2}$。允许所有分支，则得到一族正实值。两种说法都能使用，但必须说明当前约定；不能先采用多值对数，再在计算中悄悄把它当唯一函数。

幂的复合规则也因此需要谨慎。实数正底数下熟悉的某些恒等式，推广后可能因为中间重新选取分支而失败。每次消去指数与对数之前，都应检查是否丢失整圈信息。

请把推导分成三行：选对数值、乘指数、再取指数函数。这样每一步都有明确对象与约定，也能看见 Lean 验证的是其中的代数核心，而不是建立一个没有分支问题的全局复数幂系统。
